
Le projet de programmation ''Royal Dédale'' consiste à implémenter un jeu d'exploration qui se déroule dans un environnement dynamique représenté par un manoir, dont la disposition des pièces change à chaque nouvelle partie. L'objectif pour le joueur est de naviguer dans le dedale depuis l'entrée jusqu'à une salle de sortie spécifique, en découvrant la carte au fur et à mesure de sa progression.
Le manoir est composé de tuiles carrées, chacune possédant des caractéristiques spécifiques. La possibilité d'ouverture des tuiles dans certaines directions détermine les déplacements possibles entre les pièces et définit la structure labyrinthique du jeu. Au fil de l'exploration, de nouvelles tuiles sont placées, ce qui modifie la configuration du manoir afin que le joueur s’adapte à chaque partie.\\

Ce rapport détaille ainsi la conception et l'implémentation d'une architecture de programmation robuste et extensible, permettant de générer, explorer et manipuler un manoir tout en garantissant une gestion correcte des ressources et des interactions, en tenant compte des contraintes imposées par les interfaces abstraites du projet, pour enfin garantir le bon déroulement d'une partie.
